\section{Timeline Estimates for Cryptographically Relevant Quantum Computers}

\label{sec:timeline}

\subsection{Hardware Roadmaps}

We survey published roadmaps from leading quantum hardware vendors:

\begin{table}[ht]
\centering
\caption{Quantum hardware roadmaps (published projections as of 2025).}
\label{tab:roadmaps}
\begin{tabular}{lrrl}
\toprule
\textbf{Vendor} & \textbf{Qubits (2025)} & \textbf{Projected CRQC} & \textbf{Source} \\
\midrule
IBM & 1{,}121 (Condor) & 2029--2033 & IBM Quantum Roadmap \\
Google & 105 (Willow) & 2029--2035 & Google AI Blog \\
PsiQuantum & --- & 2027--2030 & Photonic approach \\
IonQ & 36 (Forte) & 2030--2035 & Trapped ion \\
QuEra & 256 (neutral atom) & 2028--2033 & Neutral atom \\
\bottomrule
\end{tabular}
\end{table}

\subsection{Logical vs. Physical Qubits}

Current quantum computers use physical qubits with error rates of $\sim$$10^{-3}$.
Breaking secp256k1 requires $\sim$$2{,}330$ \emph{logical} qubits with error rates
$< 10^{-10}$. Using surface codes with code distance 27, each logical qubit requires
$\sim$$1{,}457$ physical qubits~\cite{gidney2021factor}. Therefore:

\[
  2{,}330 \times 1{,}457 \approx 3{,}394{,}810 \text{ physical qubits}
\]

This is approximately $3{,}000\times$ more physical qubits than the largest current
systems. However, qubit counts have been growing exponentially, and error correction
improvements may reduce the overhead factor significantly.

\subsection{Probability Assessment}

Based on published roadmaps, error correction progress, and the distinction between
physical and logical qubits, we estimate:

\begin{table}[ht]
\centering
\caption{Cumulative probability of CRQC existence.}
\label{tab:crqc-probability}
\begin{tabular}{lr}
\toprule
\textbf{Year} & \textbf{P(CRQC exists)} \\
\midrule
2028 & $< 1\%$ \\
2030 & $5\text{--}10\%$ \\
2033 & $15\text{--}30\%$ \\
2035 & $30\text{--}50\%$ \\
2040 & $> 70\%$ \\
\bottomrule
\end{tabular}
\end{table}

These estimates are consistent with Mosca's inequality~\cite{mosca2018}: if the time
to deploy quantum-safe cryptography ($d$) plus the shelf life of current secrets
($l$) exceeds the time until CRQCs exist ($q$), then action is needed now. For
blockchain, $l = \infty$ (on-chain data never expires), so the condition $d < q$
must hold---meaning migration must \emph{complete} before CRQCs exist. Given that
protocol-level cryptographic migrations take 3--5 years to fully deploy, the window
for action is already closing.

%% ========================================================================
%%  6. ECONOMIC ANALYSIS
%% ========================================================================
